\documentclass{article}
\usepackage{listings}
\usepackage{xcolor}

\lstset{
	basicstyle=\ttfamily\small,
	backgroundcolor=\color{gray!10},
	frame=single,
	breaklines=true
}

\begin{document}
	
	\section*{Per-project PDF output directory in TeXstudio (Windows)}
	
	\subsection*{1. Create the batch script}
	
	Create a file named \texttt{compile.bat} in the project directory with the following content:
	
	\begin{lstlisting}[language=bash]
		@echo off
		pdflatex -synctex=1 -interaction=nonstopmode ^
		-output-directory="C:/Git/GitHub/library/library/documentation/modularity" %1
	\end{lstlisting}
	
	\subsection*{2. Configure TeXstudio}
	
	Go to \textbf{Options} $\to$ \textbf{Configure TeXstudio} $\to$ \textbf{Build} $\to$ \textbf{Commands} $\to$ \textbf{PdfLaTeX} and set:
	
	\begin{lstlisting}
		pdflatex.exe -synctex=1 -interaction=nonstopmode
		-output-directory="C:/Git/GitHub/library/library/documentation/modularity" "%"
	\end{lstlisting}
	
	\subsection*{Notes}
	
	\begin{itemize}
		\item The output directory must exist before compilation. \texttt{pdflatex} does not create it automatically.
		\item Auxiliary files (\texttt{.aux}, \texttt{.log}, \texttt{.toc}, etc.) are also written to the output directory.
		\item The setting in \textbf{Commands} applies globally to all projects. For per-project configuration, use the \texttt{.bat} script and point the \textbf{PdfLaTeX} command to it:
	\end{itemize}
	
	\begin{lstlisting}
		"C:/path/to/project/compile.bat" "%"
	\end{lstlisting}
	
\end{document}